\subsection{A logarithmic endpoint for integrable data}
\label{subsec:unibranch-curve-density-endpoint}
\index{Unibranch space!integrable density endpoint}

The hypothesis $p>1$ in \cref{thm:unibranch-torus-curve-equation} guarantees that the Green correction is bounded and continuous. For merely integrable densities, the solution can leave the ordinary type. The following family of densities locates the boundary exactly, and it also gives continuous residuals for some densities outside every $L^p$, $p>1$.

We keep the notation $X,E,h,\beta,\beta_0,H,\Psi,\psi,A,M$ of \cref{sec:unibranch-curve-equations}, with $A<H$. Fix $q\in E$ and a flat coordinate $z$ centered at $q$, and put $r=|z|$ and $t=\log(1/r)$, so that $t>1$ on a small punctured disk around $q$. For $C_0>0$ and $\gamma>1$, choose a positive density $f_\gamma$, smooth away from $q$, which equals
\[
 f_\gamma=C_0r^{-2}t^{-\gamma}
\]
on a sufficiently small punctured disk around $q$, and extend it smoothly and positively to the complement. Set
\[
 F_\gamma=\int_Ef_\gamma\beta_0,\qquad
 c_\gamma=M/F_\gamma,\qquad
 \nu_\gamma=\pi_*(f_\gamma\beta_0).
\]
The local integral of $f_\gamma$ is finite exactly for $\gamma>1$. For every such $\gamma$, this density lies outside $L^p(E,\beta_0)$ for every $p>1$. Since $f_\gamma\beta_0$ is absolutely continuous on $E$, the measure $\nu_\gamma$ is non-pluripolar by the pullback argument in the proof of \cref{thm:unibranch-torus-curve-equation}.

\Needspace{13\baselineskip}
\begin{theorem}\label{thm:unibranch-curve-density-endpoint}
For this family, the equation $\mathrm{MA}^{\mathrm{ext}}_{X,h}(u)=c\nu_\gamma$ with some constant $c>0$ has a solution in the ordinary type of $\psi$ if and only if $\gamma>2$. In that case $c=c_\gamma$, the normalized solution is unique in the whole type, and its residual is continuous but, on the normalization, not H\"older continuous for any positive exponent.

For $1<\gamma\leq2$, there is a solution $u_\gamma\in\PSH(X,h)$ of
\[
 \mathrm{MA}^{\mathrm{ext}}_{X,h}(u_\gamma)
          =c_\gamma\nu_\gamma
\]
with non-pluripolar mass $M$, whose pullback to $E$ has the same multiplier ideals as $\Psi$ at every positive weight. It violates the lower inequality $\psi\leq u_\gamma+C$ of the ordinary type, and no solution with any positive constant exists in the ordinary type of $\psi$.
\end{theorem}
\begin{proof}
The coefficient below differs from that in \eqref{eq:curve-log-primitive} because here $\beta_0=dA/V$. Define on the punctured coordinate disk
\begin{equation}\label{eq:unibranch-density-primitive}
 P_\gamma(t)=
 \begin{cases}
 \displaystyle\frac{4\pi C_0}{V(\gamma-1)(\gamma-2)}
                         t^{2-\gamma},&\gamma\ne2,\\[4pt]
 \displaystyle-\frac{4\pi C_0}{V}\log t,&\gamma=2.
 \end{cases}
\end{equation}
Direct differentiation gives
\[
 P_\gamma''(t)=\frac{4\pi C_0}{V}t^{-\gamma},\qquad
 \ddc P_\gamma=f_\gamma\beta_0
\]
off $q$. The flux of $P_\gamma$ through the circle of radius $\mathrm{e}^{-t}$, which computes the mass of $\ddc P_\gamma$ inside this circle, is
\[
 -\tfrac12P_\gamma'(t)
       =\frac{2\pi C_0}{V(\gamma-1)}t^{1-\gamma}\longrightarrow0
\]
as $t\to\infty$. Thus the distributional extension across $q$ adds no atom. Also $P_\gamma$ is locally integrable and $|P_\gamma(t)|/t\to0$, so its Lelong number at $q$ is zero. These are the same flux and growth mechanisms as in \cref{thm:curve-products}, with the present normalization of area.

Take a smooth cutoff $\chi$ supported in the coordinate disk and equal to one near $q$, and extend $L=c_\gamma\chi P_\gamma$ by zero. The signed measure
\[
 \rho\beta_0=(c_\gamma f_\gamma-M)\beta_0-\ddc L
\]
has a smooth density $\rho$ on all of $E$: near $q$ the singular densities cancel, leaving $-M\beta_0$. Its integral is zero by Stokes' theorem and $c_\gamma F_\gamma=M$. Hence \cref{lem:torus-logarithmic-green} gives a smooth correction $R=G*(\rho\beta_0)$. Set $v_\gamma=L+R$ and $U_\gamma=\Psi+v_\gamma$. Then
\begin{equation}\label{eq:unibranch-density-current}
 \beta+\ddc U_\gamma
       =c_\gamma f_\gamma\beta_0+\sum_j a_j\delta_{q_j}.
\end{equation}
\Needspace{7\baselineskip}
As before, the positive local Laplacians give subharmonic representatives. The disk-lifting argument in the proof of \cref{thm:unibranch-torus-curve-equation} descends $U_\gamma-\sup_EU_\gamma$. It is smooth away from the original poles and from $q$. Its non-pluripolar product removes the original atoms, and a possible new value $-\infty$ at $q$ carries no atom. The mass is therefore $M$, and \eqref{eq:space-extended-ma} gives the equation on $X$.

For $\gamma>2$, $P_\gamma\to0$, so $v_\gamma$ extends continuously and boundedly, and the solution belongs to the prescribed ordinary type. The difference of $v_\gamma$ from its value at $q$ has a nonzero leading term proportional to $t^{2-\gamma}$, with a smooth remainder of order $O(r)$. This decays more slowly than every power $r^\alpha$, $\alpha>0$, so the residual is not H\"older continuous for any positive exponent. Uniqueness follows from the mass of the whole type, the restoration of the fixed pole atoms and the harmonic-difference argument used in the proof of \cref{thm:unibranch-torus-curve-equation}; none of these steps requires $p>1$ once the bounded correction exists.

For $1<\gamma\leq2$, $P_\gamma\to-\infty$. The correction is bounded above but unbounded below near $q$, so the inequality $\Psi\leq U_\gamma+C$ fails off the pole set. Since its Lelong number is zero, the point atoms of $\beta+\ddc U_\gamma$ are those of $\beta+\ddc\Psi$. Applying \cref{lem:planar-multiplier-ideal} to each positive multiple, after a smooth local correction, shows that all multiplier ideals of $U_\gamma$ and $\Psi$ on $E$ coincide, also when $q$ is not one of the original poles. On $X$, this is resolution-ideal equivalence in the sense defined before \cref{prop:space-resolution-i-envelope}.

If another solution belonged to the ordinary type, its mass would be $M$, which forces the constant $c_\gamma$. Restoring its prescribed atoms in the extended equation would give exactly \eqref{eq:unibranch-density-current}. Its difference from $U_\gamma$ would then be harmonic on the torus, hence constant, with equality of the values everywhere. This contradicts the unboundedness of the correction near $q$, so no positive constant admits a solution in the ordinary type.

Finally, the integral of the density is a constant times $\int^\infty t^{-\gamma}dt$, whereas the integral of its $p$th power, $p>1$, is comparable to
\[
 \int^\infty e^{2(p-1)t}t^{-p\gamma}dt=+\infty.
\]
This proves both integrability assertions and completes the proof.
\end{proof}

The theorem is an exact criterion for this logarithmic family only; it does not characterize general $L^1$ data, and it gives no mass or uniqueness statement for the whole resolution-ideal class. On the projective cusp, all these constructions apply with the same, possibly degenerate, pulled-back background. The failure at $\gamma=2$ occurs already in the bounded ordinary type, that is, when there are no original poles.
